\documentclass[11pt,letterpaper]{article}
\usepackage[T1]{fontenc}
\usepackage[utf8]{inputenc}
\usepackage{lmodern}
\usepackage[margin=1in,headheight=14pt]{geometry}
\usepackage{amsmath,amssymb,amsthm,mathtools}
\usepackage{microtype,booktabs,longtable,array,enumitem}
\usepackage[dvipsnames]{xcolor}
\usepackage{fancyhdr}
\usepackage[colorlinks=true,linkcolor=MidnightBlue,citecolor=MidnightBlue,urlcolor=MidnightBlue,breaklinks=true]{hyperref}
\usepackage{bookmark}
\hypersetup{pdftitle={Logarithmic-integral strengthenings of the prime number theorem in iPRA},pdfauthor={Research note prepared for Vladimir Reshetnikov}}
\setlength{\emergencystretch}{3em}
\setlength{\parskip}{4pt}
\setlength{\parindent}{0pt}
\setlist{itemsep=3pt,topsep=5pt}
\pagestyle{fancy}
\fancyhf{}
\fancyhead[L]{\small Logarithmic-integral PNT in iPRA}
\fancyhead[R]{\small Explicit finite proofs}
\fancyfoot[C]{\thepage}
\numberwithin{equation}{section}
\newtheorem{theorem}{Theorem}[section]
\newtheorem{lemma}[theorem]{Lemma}
\newtheorem{corollary}[theorem]{Corollary}
\theoremstyle{definition}
\newtheorem{definition}[theorem]{Definition}
\newcommand{\ipra}{\ensuremath{\mathrm{iPRA}}}
\newcommand{\Li}{\operatorname{Li}}
\newcommand{\N}{\mathbb N}
\newcommand{\Q}{\mathbb Q}

\begin{document}
\hypersetup{pageanchor=false}
\begin{titlepage}
\vspace*{0.6in}
{\color{MidnightBlue}\Large Constructive analytic number theory\par}
\vspace{0.35in}
{\huge\bfseries Logarithmic-integral strengthenings\par
of the prime number theorem\par}
\vspace{0.18in}
{\LARGE Explicit finite proofs for intuitionistic\par
primitive recursive arithmetic\par}
\vspace{0.5in}
{\large Prepared for Vladimir Reshetnikov\par}
{\large November 15, 2026\par}
\vspace{0.45in}
\begin{abstract}
We give a quantitative constructive proof of
\[
 \pi(n)-\Li_2(n)
 =O\!\left(\frac n{\log n}\exp[-c(\log n)^{1/10}]\right)
\]
with an explicit positive constant and primitive recursive thresholds.
It follows that the error is smaller than every fixed inverse power of
$\log n$, that the full factorial asymptotic expansion holds with an
explicit remainder, and that every fixed truncation has an effectively
controlled positive first-omitted-term error. We also obtain a bounded
search procedure finding a prime in a specified relative interval.

The proof replaces nonvanishing by numerical apartness, uses rational
contours and a finite smoothed Perron kernel, and provides explicit
truncation and rounding schedules. Its main first-order statement is an
integer inequality involving a fixed rational algorithm for $\Li_2$.
The article includes the human proof, a quantifier-free induction audit,
and a detailed route to an embedded Lean proof in the ProveIt presentation
of iPRA. The included Lean sources are uncompiled outlines; no completed
mechanized PNT proof is claimed.
\end{abstract}
\vfill
\fbox{\begin{minipage}{0.92\textwidth}
\textbf{Scope.} The error exponent proved here is $1/10$. The classical
square-root-logarithm and Vinogradov--Korobov rates, and conditional
Riemann-hypothesis bounds, are discussed for comparison and are not
claimed as results of this package.
\end{minipage}}
\end{titlepage}
\hypersetup{pageanchor=true}
\pagenumbering{roman}

\begingroup
\small
\setlength{\parskip}{0pt}
\tableofcontents
\endgroup
\newpage
\pagenumbering{arabic}

\section{Results and proof strategy}

The ordinary prime number theorem is
$\pi(x)\sim x/\log x$, equivalently $\pi(x)\sim\Li_2(x)$.
An estimate for the \emph{difference} $\pi(x)-\Li_2(x)$ supplies
substantially more information. The question asks for such strengthenings
with the logical resources restricted to intuitionistic primitive recursive
arithmetic. We interpret the abbreviation ``PNM'' in the question as PNT.

The exact target is Theorem~\ref{thm:object}. In conventional notation,
write $C=2^{240}$ and $D=2C$. The following are consequences proved below:
\begin{align*}
 n\ge2^{D(m+5)^{10}}
 &\ \Longrightarrow\
 |\pi(n)-\Li_2(n)|\le2^{-m}\frac n{\log n},\\
 n\ge2^{D6^{10}}
 &\ \Longrightarrow\
 |\pi(n)-\Li_2(n)|\le
 64\frac n{\log n}\exp\!\left(-2^{-26}(\log n)^{1/10}\right).
\end{align*}
For every $k\ge0$, define
\[
 M(k)=1+6400C^2k^2+48C.
\]
Then the fully specified threshold
$n\ge2^{2M(k)^{10}}$ implies
\[
 |\pi(n)-\Li_2(n)|\le \frac n{(\log n)^{k+1}}.
\]
The constants are deliberately coarse so that their proofs and arithmetic
representations stay transparent. The exponent $1/10$ is a classical
nonoptimal exponent; this note does not claim a new classical exponent.

For $r\ge1$, put
\[
 S_r(n)=n\sum_{j=0}^{r-1}\frac{j!}{(\log n)^{j+1}},
 \qquad B_r=1+r!(1+2^{r+2}).
\]
An explicit PR threshold in $r$ gives
\[
 |\pi(n)-S_r(n)|\le B_r\frac n{(\log n)^{r+1}},
\]
and a further explicit threshold gives
\[
 \frac{r!}{2}\frac n{(\log n)^{r+1}}
 \le\pi(n)-S_r(n)
 \le\frac{3r!}{2}\frac n{(\log n)^{r+1}}.
\]
This is a statement about each fixed finite truncation. The factorial
series is an asymptotic expansion; we make no convergence assertion for
it at a fixed argument.

\begin{center}
\begin{tabular}{@{}p{0.55\textwidth}p{0.37\textwidth}@{}}
\toprule
Estimate or consequence & Status in this package\\
\midrule
Explicit stretched exponential with exponent $1/10$ & Proved, with arithmetic modulus\\
Every inverse-logarithmic order & Proved, with uniform PR thresholds\\
Factorial expansion and first omitted term & Proved, with boundary terms retained\\
Prime in a prescribed relative interval & Proved, with bounded search witness\\
Classical $e^{-c\sqrt{\log x}}$ error & Stronger benchmark; not proved here\\
Vinogradov--Korobov error & Stronger benchmark; not proved here\\
$O(\sqrt x\log x)$ and Schoenfeld's bound & Require their RH assumption; not proved here\\
\bottomrule
\end{tabular}
\end{center}

The hard step is an explicitly apart zero-free strip. For a positive
integer $Q\ge\log(|t|+3)$ and $|t|\ge1/64$, we prove
\[
 \sigma\ge1-\frac{2^{-88}}{Q^9},\quad \sigma\le2
 \quad\Longrightarrow\quad
 |\zeta(\sigma+it)|\ge2^{-71}Q^{-7}.
\]
The lower bound licenses division by $\zeta$ at a known precision.
It comes from a finite Euler-product inequality and a horizontal
derivative bound, so no search for zeros or compactness argument is used.

Next, integrate the prime-power counting function once. The kernel
$1/(s(s+1))$ gives an absolutely convergent vertical integral with a
numerical tail bound. A finite rectangle calculation moves the contour
past the single explicitly isolated pole at $s=1$. Monotonicity then
removes the smoothing. Finally, exact finite Abel summation transfers
the resulting estimate from $\vartheta$ to $\pi-\Li_2$.

The conventional analytic method is familiar from classical accounts
\cite{Kontorovich,Tao}. The work required by the present setting lies
in its lower bounds, finite precisions, and quantifier-free induction
invariants. These are included as part of the proof, rather than delegated
to a general claim about computable analysis.

\section{The precise arithmetic meaning of the result}
\label{sec:foundations}

We use first-order intuitionistic primitive recursive arithmetic, abbreviated
\ipra. Its language contains a symbol for every primitive recursive function,
with its defining equations; its logic is intuitionistic first-order logic
with equality, and its induction schema is restricted to quantifier-free
formulas. Natural numbers are the only objects. Integers, positive-denominator
rationals, Gaussian rationals, and finite lists are coded by natural numbers.
No quantification over a completed space of real or complex numbers is added.

The corresponding presentation in the inspected repository is
\texttt{PRAAx} in \texttt{ConciseIPRA/PR.lean}; despite its identifier, the
ambient deduction system is intuitionistic. Its induction interface is
\texttt{PRA.qfInd}. The theorem \texttt{pra\_iff\_native} transports a
\emph{derivation already constructed in this theory} to the native scalar
calculus. It does not turn an arbitrary theorem of ambient Lean into a
theorem of \ipra. All repository references in this package are pinned to
commit \texttt{b5a3ad037723c44d69a3d5b28a950c554327a989}; see
\cite{ProveIt} and the evidence manifest.

\subsection{Fixed algorithms, rather than a universal real-number sort}

An expression such as $\log u$, $e^z$, or $\zeta(s)$ below abbreviates
one specified family of rational computations, together with an error
bound. A finite assertion with requested error $2^{-v}$ is proved using
the finite computation at a prescribed finer precision. The notation for
the limiting value helps the reader see the mathematics; the arithmetic
claim is about those finite enclosures. In particular, we never postulate
a primitive recursive evaluator for all primitive recursive program codes.
The small, fixed collection of algorithms used here is compiled separately.

Bounded quantifiers cause no increase of the induction schema. For example,
if $R(i,\mathbf a)$ is a decidable primitive recursive predicate, then
``$R(i,\mathbf a)$ for all $i< B$'' is represented by the equation
\[
 \prod_{i<B}\chi_R(i,\mathbf a)=1.
\]
Products, sums, maximum, finite search with a supplied upper bound, and
the characteristic function in this equation are primitive recursive.
Induction applies to this single equation with all other parameters free.
The proof may universally generalize those parameters afterwards.

\subsection{A concrete logarithm algorithm}

For a positive rational $u\ge1$, find, by bounded integer search, the unique
$e\ge0$ with $2^e\le u<2^{e+1}$ and put $r=u/2^e$. Define
\[
 A_j(r)=2\sum_{i=0}^{j-1}\frac{z^{2i+1}}{2i+1},
 \qquad z=\frac{r-1}{r+1},\qquad 1\le r\le2.
\]
The finite geometric identity, integrated on $[0,z]$, gives
\[
 0\le \log r-A_j(r)
 =2\int_0^z\frac{t^{2j}}{1-t^2}\,dt
 \le\frac{2z^{2j+1}}{(2j+1)(1-z^2)}\le9^{-j}.
\]
This also proves convergence with an explicit bound, without invoking a
monotone-convergence principle. Set
\begin{equation}\label{eq:log-algorithm}
 \lambda(u,v)=eA_{v+e+2}(2)+A_{v+e+2}(r).
\end{equation}
Then
\begin{equation}\label{eq:log-error}
 0\le\log u-\lambda(u,v)\le(e+1)9^{-(v+e+2)}\le2^{-v}.
\end{equation}
The last inequality follows from $e+1\le2^{e+1}$ and $9\ge 2^3$.
For $u\ge2$, the initial term of $A_j(2)$ also gives
$\lambda(u,v)\ge2/3$. Rational inputs below $1$ are handled by
$\log u=-\log(1/u)$. Every loop has a bound supplied by the numerator,
denominator, and requested precision.

The ordinary logarithm identities used later can be proved directly for
these computations. One convenient common reference is
$\log u=\int_1^u dt/t$: the geometric identity above evaluates this
integral, substitution in finite Riemann sums proves
$\log(ab)=\log a+\log b$, and the derivative and monotonicity follow from
the explicit rational identity for $1/(u+h)-1/u$. The needed error bounds
on an interval with $u\ge a>0$ use only $a^{-1}$ and $a^{-2}$.
This gives finite rational proofs of these identities at every prescribed
precision; it does not appeal to an unspecified construction of real analysis.

\subsection{An explicit rational approximation to the offset integral}

We use the convention
\[
 \Li_2(x)=\int_2^x\frac{dt}{\log t}.
\]
It has no singular endpoint. Let $b(n)$ be the binary length of the
integer $n\ge2$, so $2^{b(n)-1}\le n<2^{b(n)}$. For $v\ge0$ set
\[
 h=2^{-(v+3)},\quad J=(n-2)2^{v+3},\quad
 P=v+b(n)+6,
\]
and define the nonnegative rational number
\begin{equation}\label{eq:Li-algorithm}
 \mathcal A(n,v)=h\sum_{j=0}^{J-1}
       \frac1{\lambda(2+jh,P)}.
\end{equation}
An empty sum, in particular at $n=2$, is zero. For $n<2$ we define the
algorithm to return zero, thereby making it total on all natural inputs.

Here is a complete error estimate. The function $f(t)=1/\log t$ is
positive and decreasing on $[2,n]$. The left and right Riemann sums on
this uniform mesh differ by $h(f(2)-f(n))\le2h$. Thus the exact left
sum exceeds the integral by at most $2h$. Since both logarithm
denominators are at least $2/3$, replacing $\log t$ by its lower
approximation increases its reciprocal by at most $3\cdot2^{-P}$.
Summing this error and using $n<2^{b(n)}$ gives
\begin{equation}\label{eq:Li-error}
 0\le\mathcal A(n,v)-\Li_2(n)
 \le 2h+3n2^{-P}
 <\left(\frac14+\frac3{64}\right)2^{-v}<2^{-v}.
\end{equation}
The executable reference implementation in the archive computes this
same rational sum. It is intended to specify a transparent algorithm,
not to provide an efficient way to evaluate the integral for very large $n$.

\subsection{The main first-order theorem}

Write $\Pi(n)$ for the number of integers $p\le n$ satisfying
$p\ge2$ and having no proper divisor between $2$ and $p-1$.
This definition uses only bounded divisibility tests and finite summation.
Set
\[
 C=2^{240},\qquad D=2C=2^{241},\qquad
 \mathcal N(m)=2^{D(m+7)^{10}}.
\]
The principal arithmetic conclusion of the argument is
\begin{theorem}[Finite logarithmic-integral strengthening]\label{thm:object}
In \ipra, for all natural $m,v,n$,
\begin{equation}\label{eq:object-target}
 n\ge\mathcal N(m)\quad\Longrightarrow\quad
 \left|\Pi(n)-\mathcal A(n,v)\right|
 \le \frac{n}{2^m b(n)}+2^{-v}.
\end{equation}
All functions and comparisons in this statement are primitive recursive.
After clearing the positive denominators, its matrix is quantifier-free.
\end{theorem}

For example, write $\mathcal A=a/d$ with $d>0$, without requiring the
fraction to be reduced, and let $z=|d\Pi(n)-a|$. The consequent is exactly
the integer comparison
\begin{equation}\label{eq:integer-matrix}
 z\,2^{m+v}b(n)
 \le d\bigl(n2^v+2^m b(n)\bigr).
\end{equation}
Thus the displayed theorem does not quantify over real numbers and does
not use an undecidable comparison between real-number codes.

The constants are chosen for simple verifiable estimates. They are
enormously larger than constants useful for numerical prime counting.
No claim of numerical optimization or a new classical error exponent is made.

\subsection{Elementary complex computations and their resource bounds}

For a Gaussian rational $z$ with $|z|\le B$, evaluate the exponential by
its Taylor polynomial. If $J$ is even and
$J\ge\max\{4\lceil B\rceil^2,2v+4\}$, then
\[
 \left|e^z-\sum_{j<J}\frac{z^j}{j!}\right|\le2^{-v}.
\]
Indeed, $J!\ge(J/2)^{J/2}$ gives $B^J/J!\le2^{-J/2}$.
Successive omitted terms then decrease by a factor at most $1/2$,
so their sum is at most $2^{1-J/2}\le2^{-v}$.
We always use the unambiguous safe choice
\begin{equation}\label{eq:exp-cutoff}
 J=4\lceil B\rceil^2+2v+4.
\end{equation}
Sine and cosine are the imaginary and real parts at $iz$; powers
$u^{-s}$ are evaluated as $e^{-s\log u}$. When the logarithm is also
approximated, a bound on the exponential derivative prescribes its
input precision. For $|z|\le B$ one may use the rational bound
$|e^z|\le3^{\lceil B\rceil+1}$, and ask for an input error at most
$2^{-v}/(3^{\lceil B\rceil+2}+1)$. Bit lengths turn this into a
natural-number precision by a bounded calculation.

Division is used only with a displayed lower bound. If $|w|\ge a>0$
and $|w-\widetilde w|\le\eta\le a/2$, then
\[
 |\widetilde w|\ge a/2,
 \qquad |w^{-1}-\widetilde w^{-1}|\le2\eta/a^2.
\]
The same estimate is obtained using squared norms and rational
arithmetic; square roots need not be decided. The zeta lower bounds
proved below are exactly what makes this operation legitimate.

For a real-parameter integral on a rational interval of length $\ell$,
a derivative bound $M$ makes a uniform mesh of width $h$ incur an error
at most $M\ell h$. A finite sum of $J$ terms is evaluated with
per-term error at most $2^{-v}/(J+1)$, adjusted for rational coefficients.
Products use their explicit magnitude bounds. These prescriptions give
precision bounds by addition, multiplication, exponentiation, and binary
length; no search for a successful accuracy is required.

There is a useful explicit rule for composing these estimates. Consider
the bounded table of $J$ rational arithmetic gates used in one calculation.
Let an integer $B\ge1$ bound the magnitudes of all exact intermediate
values and let every inverted denominator have modulus at least $1/B$.
Addition, multiplication, absolute value, and guarded inversion amplify
an input error at most $\epsilon\le1/(2B)$ by at most $8B^3$, including
one further local rounding error. Set the initial tolerance to
\[
 \epsilon_0=\frac{2^{-p}}{2B(8B^3)^{J+1}}.
\]
Quantifier-free induction on the number of evaluated rows proves that
the row-$j$ error is at most $(8B^3)^j\epsilon_0$ and that every guard
remains valid. The final error is less than $2^{-p}$. The bounds $J,B$
come from the displayed cutoffs and magnitude estimates. This is an
error-propagation lemma for a bounded arithmetic table with the specified
analytic macros, not a universal evaluator for primitive recursive codes.

\subsection{Why the proof uses only quantifier-free induction}

Every analytic operation used in the number-theoretic proof is reduced
below to these operations and two explicit tail estimates. The finite
contour lemma uses a predetermined rectangular grid, and its proof is
exact cancellation of internal edges plus a rational error sum. Euler
products are cut off by a supplied integer bound. The Perron integral
has a numerical truncation bound and an explicit finite scalar-kernel
calculation. No enumeration of zeta zeros, compactness of a continuum,
unbounded minimization, or principle of choice occurs.

Here is the induction discipline, which is essential to the claim of
Theorem~\ref{thm:object}. Fix the input numbers and the final accuracy.
Compute all resource bounds first. For a loop with counter $j\le J$,
the induction invariant is a finite list of rational identities and
inequalities for the first $j$ terms, with their accumulated error
allowances. Replace that list by the conjunction of its decidable
characteristic equations. Induct on this quantifier-free formula.
Induction never asserts ``the identity holds at every precision'' or
``every point of an infinite contour has the desired property''.
A claimed uniform analytic bound is instantiated at the finitely many
mesh points and at the finitely many interval bounds required by the
Taylor estimate.

For rational statements with positive slack, all apparent passages to
a limit are eliminated by choosing one sufficiently fine approximation
from the displayed bounds. For an exact identity in the expository
notation, use the finite version with an arbitrary prescribed error.
At the last step the spare margin in the prime-counting estimate and
\eqref{eq:Li-error} absorb these errors. Universal generalization over
the original inputs is made only after these finite proofs. This is a
description of the specific proof translation used here, not the false
general principle that every computable inequality is provable in \ipra.

\paragraph{Mechanization status.}
Theorem~\ref{thm:object} is established here at the level requested: a
human proof with its finite arithmetic interpretation. The archive does
\emph{not} contain a completed Lean proof term for it. The Lean files are
explicitly marked interfaces and outlines, and the final section lists
the proof objects still to be constructed. In particular, a successful
test of a rational evaluation program would not establish the universal
number-theoretic theorem.

% Complete quantitative analytic core; included by the main manuscript.
\section{Finite analytic constructions used below}
The complex notation in this section abbreviates pairs of rational
approximation procedures with certified errors. It does not quantify over
arbitrary real numbers or arbitrary functions. All final inputs $x,y$ are
positive rationals, all vertical contour sides used in the final estimate
are rational, and every infinite expression below is accompanied by a
tail bound. The proof only needs finite sums, products, powers, rational
division with a supplied positive denominator bound, and rational
approximations to the elementary exponential, sine, cosine and logarithm.
The latter can be obtained from their power series after range reduction.
For example, $\log u$ on a positive rational interval follows from
$\log u=2\sum_{j\ge0}z^{2j+1}/(2j+1)$, $z=(u-1)/(u+1)$;
binary range reduction puts $|z|\le1/3$, and the geometric tail is
explicit. Thus every expression used here has a primitive recursive
approximation at a requested precision.

\paragraph{Finite Euler summation.}
Let $s=\sigma+it$, $\sigma>0$, and let $1\le N<M$ be integers.
Integration by parts on each interval $[n,n+1]$ and telescoping gives
\[
 \sum_{n\le M}n^{-s}=\sum_{n\le N}n^{-s}
 +\frac{N^{1-s}-M^{1-s}}{s-1}
 -s\int_N^M\{u\}u^{-s-1}\,du.                              \tag{E1}
\]
For $s=1$ the quotient means its explicit logarithmic limit; later we
separate this point algebraically. One can verify (E1) without integration
notation: on $[n,n+1]$, replace $\{u\}$ by $u-n$, use the elementary
antiderivatives of $u^{-s}$ and $u^{-s-1}$, and sum. Both endpoint sums
telescope. The displayed integrals are therefore finite elementary
expressions before any limiting operation.

For $\sigma>1$ the term $M^{1-s}$ and the tail integral have explicit
limits, so (E1) agrees with the Dirichlet series definition of $\zeta$.
For $\sigma>0$, $s\ne1$, define instead
\[
 Z(s)=\sum_{n\le N}n^{-s}+\frac{N^{1-s}}{s-1}
       -s\int_N^\infty\{u\}u^{-s-1}\,du.                   \tag{E2}
\]
Equation (E1) proves that this does not depend on $N$. The tail after $V$
is at most $|s|V^{-\sigma}/\sigma$. This is the only continuation of
$\zeta$ used in this proof; no analytic-continuation uniqueness theorem is
needed. Differentiating the finite expression and its explicitly bounded
tails is legitimate, since for $j\ge0$
\[
 \int_V^\infty (\log u)^j u^{-\sigma-1}\,du
 =V^{-\sigma}\sum_{r=0}^j
     \frac{j!}{(j-r)!}\frac{(\log V)^{j-r}}{\sigma^{r+1}}.    \tag{E3}
\]
This identity follows by $j$ integrations by parts; in the application
$j\le3$ and $\sigma\ge1/2$. Taking $V=2^v$ with a sufficiently large
explicit integer $v$ makes each tail as small as requested. For instance,
one may start with $v=2^{q+10}(T+4)^4$ when $|t|\le T$ and the desired
error is $2^{-q}$; the exponential decay $2^{-v/2}$ dominates the degree
three polynomial in $v$ appearing in (E3). Every such comparison reduces
to the elementary inequalities $n\le2^n$ and finite multiplication.

Writing $\zeta$ for $Z$, the $N=1$ formula is
\[
 \zeta(s)=\frac1{s-1}+h(s),\qquad
 h(s)=1-s\int_1^\infty\{u\}u^{-s-1}\,du.                   \tag{E4}
\]
Formula (E3) proves directly that $h$ has three controlled complex
derivatives on every rectangle used below. The pole at $1$ is therefore
isolated by an explicit algebraic expression, rather than detected by a
zero or pole enumeration.

\paragraph{Finite Euler products and their tails.}
For $\sigma=1+\eta>1$ put
$P_N(s)=\prod_{p\le N}(1-p^{-s})^{-1}$. Every factor has a supplied
apartness bound $|1-p^{-s}|\ge1-p^{-\sigma}\ge1/2$.
Expand each factor as a geometric series. Unique prime factorization
identifies the resulting product with the sum of $n^{-s}$ over integers
whose prime factors are at most $N$. This assertion can be checked first
with all exponents bounded and then with geometric tails; no infinite
choice is involved. Every integer $n\le N$ has all its prime factors at
most $N$, and consequently
\[
        |\zeta(s)-P_N(s)|\le \sum_{n>N}n^{-\sigma}
                    \le N^{-\eta}/\eta.                   \tag{E5}
\]
The last inequality is comparison with $\int_N^\infty u^{-1-\eta}du$.
If $\eta\ge1/B$, the explicit choice
$N=2^{B(q+B+2)}$ makes the right-hand side at most $2^{-q-2}$:
$N^{-\eta}/\eta\le B2^{-q-B-2}\le2^{-q-2}$.
For a finite verification of the geometric expansion, truncate each prime
exponent at $J\ge q+2N+10$. The product error is bounded by
$N2^N2^{-J}$, because there are at most $N$ factors, each complete factor
has modulus at most $2$, and each geometric tail is at most $2^{-J}$.
This supplies a literal finite expansion with a primitive recursive size.

The elementary logarithm identity
\[
 -\log|1-z|=\mathop{\rm Re}\sum_{j\ge1}z^j/j,
                  \qquad |z|\le1/2,
\]
follows by integrating the finite geometric series for
$z/(1-uz)$ from $u=0$ to $1$. Its tail after $J$ has modulus at most
$2|z|^{J+1}/(J+1)\le2^{-J}$. Apply it to
$z=p^{-\sigma}\exp(-it\log p)$ and add with coefficients $3,4,1$ at
ordinates $0,t,2t$. The summand is
\[
 \frac{p^{-j\sigma}}j
 [3+4\cos(jt\log p)+\cos(2jt\log p)]
 =\frac{2p^{-j\sigma}}j[1+\cos(jt\log p)]^2\ge0.
\]
For the product over $p\le N$, truncation at
$J\ge q+\lceil\log_2N\rceil+10$ has total logarithmic error below
$2^{-q-7}$. The identity $\exp(u+v)=\exp(u)\exp(v)$, proved from its
absolutely convergent power series, now yields the positive product
inequality used below.

One does not have to infer apartness from an exact limiting inequality.
For a direct finite version sufficient for the constants below, suppose
$0<\eta\le1$, choose $B$ with $\eta\ge1/B$, put $q=20+8B$, and use (E5).
Then each of $\zeta(\sigma)$, $\zeta(\sigma+it)$ and
$\zeta(\sigma+2it)$ differs from its finite Euler product by at most
$\varepsilon=2^{-20}B^{-8}$. All six quantities have modulus at most
$4/\eta$. The change in a product of the eight factors, counted with
multiplicity, is at most
\[
 8\varepsilon(4/\eta)^7
 \le2^{17}2^{-20}B^{-8}\eta^{-7}\le1/8.
\]
The finite logarithmic calculation gives its Euler-product side at least
$7/8$ after allowing another explicit $1/8$ approximation error. Hence
the product of the zeta moduli is at least $3/4$, and in particular at
least $1/2$. This weaker numerical version already proves the rounded
$2^{-4}$ lower bound used below. Thus the lower bound has a finite
positive certificate; the argument does not use a principle converting
double-negated nonvanishing into apartness.

\paragraph{The logarithmic derivative series.}
Differentiate the finite Euler product and expand its geometric factors:
\[
 -P_N'(s)/P_N(s)
 =\sum_{p\le N}\frac{(\log p)p^{-s}}{1-p^{-s}}
 =\sum_{p\le N}\sum_{j\ge1}(\log p)p^{-js}.
\]
The unique factorization identification of prime powers gives
\[
 F(s):=-\zeta'(s)/\zeta(s)
              =\sum_{n\ge1}\Lambda(n)n^{-s},\quad\sigma>1. \tag{E6}
\]
This differentiation and limiting step is quantitative. The difference
$\zeta-P_N$ has the omitted-integer expansion used to prove (E5), and
the derivative of that expansion is controlled by the majorant
$\sum_{n>N}(\log n)n^{-\sigma}$; this does not differentiate an
inequality. The reciprocal of $\zeta$ on
$\sigma>1$ has the finite Euler-product bound
$|1/\zeta(s)|\le\prod_p(1+p^{-\sigma})\le\zeta(\sigma)$.
For $c=1+1/b$, $b\ge1$, comparison on each interval $[n-1,n]$ yields
\[
 \sum_{n\ge2}(\log n)n^{-c}
 \le\int_1^\infty\log(2u)u^{-c}du
 =b^2+b\log2\le2b^2.                                     \tag{E7}
\]
The tail after $N$ is at most
$N^{-1/b}[b\log(2N)+b^2]$. For any requested $2^{-q}$, the explicit
choice $N=2^{b^2(q+b+10)^2}$ makes this tail smaller than $2^{-q}$.
To check this, bound the logarithm of the polynomial factor by
$2q+6b+23$ using $\log_2r\le r$, and compare it with
$b(q+b+10)^2$. Enlarging $q$ by the binary length of a preceding
multiplier supplies any desired error budget after multiplication.

\paragraph{A finite scalar Perron kernel.}
For positive rational $v$ and $c\in[1,2]$, define the integral over a
finite vertical segment in the usual parametrized sense. Its integrand
$v^s/(s(s+1))$ has scalar poles at $0,-1$, with residues respectively
$1,-v^{-1}$. We claim
\[
 \frac1{2\pi i}\int_{c-i\infty}^{c+i\infty}
            \frac{v^s}{s(s+1)}ds=(1-v^{-1})_+.              \tag{E8}
\]
Here is a complete finite contour verification. If $v\ge1$, close the
segment on the line $\Re s=-M$, $M\ge2$. Its left edge has absolute
contribution at most $1/(2(M-1))$, because $v^{-M}\le1$ and
\[
 \int_{-\infty}^{\infty}
 \frac{dt}{(M-1)^2+t^2}=\frac\pi{M-1}.
\]
The two horizontal edges at height $U$ have combined contribution at
most $(M+2)\max(1,v^2)/U^2$, using $|s(s+1)|\ge U^2$.
The closed rectangle contains exactly the two displayed scalar poles,
so its integral is $1-v^{-1}$. If $0<v<1$, close instead on
$\Re s=M$, $M>2$. There are no poles; the right-edge and horizontal
bounds are no larger, since $v^M\le1$ and $v^\sigma\le1$ there.
At $v=1$, either closure gives zero, so no apartness from $1$ is needed.
The decision $v\ge1$ is an exact rational comparison.

For a requested error $2^{-q}$ and a rational bound $v\le Y$, the explicit
choices
\[
 M=2^{q+6},\qquad U=2^{q+6}(M+3)(Y+1)^2
\]
make all closing-edge errors below $2^{-q-3}$. The omitted tails on the
original line are bounded by $v^c/(\pi U)$ and satisfy the same bound.
Every contour here is finite and rational. The residue calculation can
be expressed by subtracting the two rational poles and applying the
finite rectangle lemma; the integral of $1/(s-a)$ is evaluated directly
on the four sides by logarithms and arctangents. Alternatively, the same
calculation follows from integrating its elementary derivative away
from $a$ on four rational subrectangles. In either case only the two
explicit real poles occur, and their winding is fixed by the inequalities
$-M<-1<0<c$.

Multiplying (E8) by $y\Lambda(n)$ with $v=y/n$, summing, and using
(E6)--(E7) proves the smoothed Perron formula below. To make every step
finite at accuracy $2^{-q}$, first truncate the Dirichlet series using
the cutoff after (E7), with $q$ enlarged to absorb $3y^2$; then use
(E8) at accuracy $2^{-q}/(1+yN\log N)$ for each of the at most $N$
terms. The order matters: first choose $N$, then the kernel accuracy,
then its auxiliary rectangle sizes. There is no circular choice of
cutoffs. Absolute values and the majorants above account for both
discarded tails. The later PNT contour height $T$ is an independent,
much smaller explicit number.

\section{An explicitly apart zero-free strip}
Put $H=|t|+3$, $L=\log H$, and $a=2^{-88}$. For
$|t|\ge 1/64$, $1-1/(2L)\le\sigma\le2$, Euler summation gives
\[
 \zeta(s)=\sum_{n\le N}n^{-s}+\frac{N^{1-s}}{s-1}
 -s\int_N^\infty\{u\}u^{-s-1}\,du,\qquad N=\lceil H^2\rceil.
\]
Here $L>1$, $\sigma>1/2$, $\log N\le3L$, $N^{1-\sigma}\le4$,
and $|s|\le H$. Consequently
\[
 |\zeta(s)|\le 2^9L,\qquad |\zeta'(s)|\le2^{15}L^2.                 \tag{A1}
\]
For completeness, the finite sum in the first estimate is at most $13L$;
the pole term is at most $256$; the integral term is at most $8/H<3$.
For the derivative, the three contributions are respectively at most
$39L^2$, $768L+16384$, and $8/N+40L/H$. Their sum is below $2^{15}L^2$.
The formula and its derivatives have effective tails: after truncating the
integral at $V$, the zeroth tail is at most $|s|V^{-\sigma}/\sigma$;
the first derivative tail is bounded by
$V^{-\sigma}/\sigma+|s|V^{-\sigma}(\log V/\sigma+1/\sigma^2)$.
The analogous $j$th derivatives have the explicit bound obtained by
integrating $u^{-\sigma-1}(\log u)^j$ by parts $j$ times.
For the bound $N^{1-\sigma}\le4$, use $N\le(10/9)H^2$, hence
$\log N\le(5/2)L$, and $e^{5/4}<4$ (since $e<3$ and $3^5<4^4$).

For $1<\sigma\le2$, the Euler product and
$3+4\cos v+\cos(2v)=2(1+\cos v)^2\ge0$ give
\[
 \zeta(\sigma)^3|\zeta(\sigma+it)|^4
 |\zeta(\sigma+2it)|\ge1.                                      \tag{A2}
\]
This identity requires no choice of a complex logarithm: apply the real
logarithm to the moduli of finite Euler factors, expand their absolutely
convergent real series, and use the displayed nonnegative summands. Pass to
finite approximations with the explicit $\sum_{n>N}n^{-\sigma}$ tail.
Using $\zeta(\sigma)\le2/(\sigma-1)$ and
$|\zeta(\sigma+2it)|\le 2^{10}L$ gives
\[
 |\zeta(\sigma+it)|\ge2^{-4}(\sigma-1)^{3/4}L^{-1/4}.
\]
Set $\delta=aL^{-9}$ and $\sigma_0=1+\delta$.
At $\sigma_0$ this lower bound is $2^{-70}L^{-7}$.
Moving horizontally from $\sigma_0$ by at most $2\delta$ changes $\zeta$
by at most $2^{16}aL^{-7}=2^{-72}L^{-7}$, by (A1).
For $\sigma\ge\sigma_0$ the direct product bound is at least as strong.
Thus, on the whole region
\[
 1-2^{-88}L^{-9}\le\sigma\le2,\quad |t|\ge1/64,
\]
we have the positive, numerical apartness estimates
\[
 |\zeta(s)|\ge2^{-71}L^{-7},\qquad
 \left|\frac{\zeta'(s)}{\zeta(s)}\right|\le2^{86}L^9.             \tag{A3}
\]
In particular this is stronger constructively than merely asserting that
$\zeta(s)\ne0$.

For the finite arithmetization, use the following rational-parameter
version of exactly the same argument. Let $Q$ be a positive integer with
$Q\ge\log(|t|+3)$; take $\delta=2^{-88}Q^{-9}$ and the rational base
point $\sigma_0=1+\delta$. Estimates (A1) hold with $Q$ in place of $L$
on $\sigma\ge1-1/(2Q)$. At $\sigma_0$ the product inequality, even in
its finite $1/2$ version above, gives
$|\zeta(\sigma_0+it)|\ge2^{-4}\delta^{3/4}Q^{-1/4}
=2^{-70}Q^{-7}$. The travel loss is at most $2^{-72}Q^{-7}$.
Thus, throughout $1-2^{-88}Q^{-9}\le\sigma\le2$,
\[
 |\zeta(\sigma+it)|\ge2^{-71}Q^{-7},\qquad
 |\zeta'/\zeta(\sigma+it)|\le2^{86}Q^9.                    \tag{A3r}
\]
The fractional powers in this explanation can be eliminated by fourth
powers: the product lower bound is at least
$2^{-14}\delta^3/Q$, whereas $(2^{-4}\delta^{3/4}Q^{-1/4})^4
=2^{-16}\delta^3/Q$. At the chosen $\delta$, the latter is the fourth
power of the rational number $2^{-70}Q^{-7}$. No irrational sample
point or root-selection operation enters the numerical proof.

\paragraph{The pole neighbourhood.}
Write $\zeta(s)=1/(s-1)+h(s)$, where
$h(s)=1-s\int_1^\infty\{u\}u^{-s-1}\,du$.
If $|\sigma-1|\le1/64$ and $|t|\le1/64$, then $|h|\le4$, $|h'|\le10$,
and $|s-1|\le1/32$. Consequently
$w(s)=(s-1)\zeta(s)=1+(s-1)h(s)$ satisfies $|w|\ge7/8$ and
$|w'|\le5$. Thus
\[
 -\zeta'(s)/\zeta(s)=1/(s-1)+g(s),\qquad |g(s)|\le10,            \tag{A4}
\]
with explicitly computable holomorphic $g=-w'/w$.
For an integer $k\ge2$, set $T=2^k$, $Q=k+2$, $\delta_T=2^{-88}Q^{-9}$,
$\alpha=1-\delta_T$. On the vertical edge $s=\alpha+it$, $|t|\le T$,
(A3)--(A4) and $\log(T+3)\le Q$ imply $|\zeta'/\zeta|\le2^{90}Q^9$.
On the top and bottom horizontal edges, $|\zeta'/\zeta|\le2^{86}Q^9$.
There is exactly the explicitly isolated pole at $s=1$ in the rectangle
$\alpha\le\sigma\le1+1/64$, $|t|\le T$; no zero enumeration is used.

\section{Smoothed Perron and error estimates}
For rational $y\ge e^{64}$ put $R=\log y$. Let $b$ be the least integer
with $y\le2^b$, and put $c=1+1/b$. Then $R\le b\le3R$,
$c\le1+1/64$, and $y^{c+1}\le3y^2$. Define
\[
 \Psi_1(y)=\sum_{n\le y}(y-n)\Lambda(n),\qquad F=-\zeta'/\zeta.
\]
The one-smoothed Perron identity is
\[
 \Psi_1(y)=\frac1{2\pi i}\int_{c-i\infty}^{c+i\infty}
 F(s)\frac{y^{s+1}}{s(s+1)}\,ds.                              \tag{A5}
\]
To prove it, expand $F(s)=\sum\Lambda(n)n^{-s}$ using the absolutely
convergent Euler product and evaluate
$\frac1{2\pi i}\int_{c-i\infty}^{c+i\infty}v^s/(s(s+1))\,ds$
as $(1-v^{-1})_+$, by closing left or right. The only scalar poles are
$0,-1$. At $v=1$ close right; its integral tends to zero like $1/M$ on
$\Re s=M$. When $v\ne1$ is rational, the sign of $v-1$ is decidable and
its distance from $1$ has an explicit rational lower bound. Therefore the
auxiliary rectangle sizes can all be chosen by primitive recursion.
All interchanges are justified by the numerical majorant
$\sum\log(n)n^{-c}\le2b^2\le18R^2$ and $\int(1+t^2)^{-1}dt<4$.

Truncating (A5) at $T$ costs at most $54y^2R^2/T$.
Shift the finite contour to $\alpha=1-\delta_T$. The residue is $y^2/2$.
Since $\int_{-T}^T|s(s+1)|^{-1}dt\le\pi/\alpha\le2\pi$,
the vertical edge costs at most $2^{91}y^2Q^9e^{-\delta_TR}$.
The two horizontal edges cost at most $2^{88}y^2Q^9/T^2$.
Hence
\[
 \left|\Psi_1(y)-\frac{y^2}2\right|
 \le y^2\left(2^{91}Q^9e^{-\delta_TR}
      +2^{88}Q^9/T^2+54R^2/T\right).                       \tag{A6}
\]
The contour step uses the strictly finite rectangle identity proved in Section~\ref{sec:finite-contour-audit}.
Subtract $y^2/(2(s-1))$ from the integrand. In the pole neighbourhood,
write its apparent divided difference as
$\int_0^1 A'(1+u(s-1))du$, with $A(s)=y^{s+1}/(s(s+1))$.
Away from that neighbourhood use (A3). The resulting analytic function
and its first two derivatives have the explicit integer bound
\[
 M=2^{512}(\lceil y\rceil+1)^8(\lceil T\rceil+4)^{40}.
\]
Indeed $|h^{(j)}|\le2^{10}(T+3)$ for $j\le3$, by
$\int_1^\infty(\log u)^ju^{-\sigma-1}du=j!\sigma^{-j-1}$.
For $|t|\ge1/64$, the zeta derivatives through order $3$ are at most
$2^{28}(T+3)$, its reciprocal is at most $2^{71}(T+3)^7$, and the
logarithmic derivative and its first two derivatives are at most
$2^{400}(T+3)^{32}$. For $|t|\le1/64$, use $|w|\ge7/8$ and the
integral divided-difference formula above; derivatives through order $3$
of $A(s)$ are at most $2^{12}y^6$. These bounds are all below $M$.
All contour corners are rational. A rectangular
mesh of side at most $2^{-q}/(100(M+1)(T+1)^2)$, where $M$ bounds the
second derivative, makes the sum of cell-boundary errors below $2^{-q}$.
Each cell uses the first-order Taylor polynomial, whose contour integral
is exactly zero; the $O(M\,\text{side}^3)$ errors sum to the asserted bound.
This is finite telescoping and does not invoke a general compactness theorem.

\section{A convenient all-integer scale}
Let $D_0=2^{100}$ and suppose $k\ge8192$ and
\[
 D_0 k^{10}\le\log_2x<D_0(k+1)^{10}.
\]
Choose $T=2^k$ and use the same $T$ for $y\in[x/2,2x]$.
Then $Q\le2k$, $\delta_T\ge2^{-97}k^{-9}$, and
\[
 \log y\ge2^{98}k^{10},\qquad \log y\le2^{111}k^{10}.
\]
Thus $\delta_T\log y\ge2k$.
For $k\ge8192$, $\log_2k\le k/256$. Here is a wholly integer proof.
Put $j=\lfloor k/256\rfloor\ge32$. At $j=32$ one has
$256(j+1)\le2^j$; the same inequality follows for $j+1$ from
$256(j+2)\le2\cdot256(j+1)$. Quantifier-free induction therefore gives
$k<256(j+1)\le2^j$, whence $k^{256}\le2^{256j}\le2^k$.
Thus the needed polynomial absorptions can be proved without real
differentiation: their $256$th powers are bounded respectively by
$2^{25600+9k}\le2^{32k}$ and
$2^{58368+20k}\le2^{64k}$. Consequently
$2^{100}k^9\le2^{k/8}$ and $54\cdot2^{222}k^{20}\le2^{k/4}$.
The three terms in (A6) are bounded respectively by
$2^{-15k/8}$, $2^{-15k/8}$, and $2^{-3k/4}$.
Their sum is at most $2^{-k/2}$.
Consequently, uniformly for $y\in[x/2,2x]$,
\[
 \left|\Psi_1(y)-y^2/2\right|\le y^2 2^{-k/2}
 \le4x^2 2^{-k/2}.                                         \tag{A7}
\]
Take $h=x2^{-\lfloor k/4\rfloor}$. Since $\psi$ is nondecreasing,
\[
 \frac{\Psi_1(x)-\Psi_1(x-h)}h\le\psi(x)
 \le\frac{\Psi_1(x+h)-\Psi_1(x)}h.
\]
It follows from (A7) that
$|\psi(x)-x|\le h/2+8x^2 2^{-k/2}/h\le9x2^{-k/4}$.
Since $9\le2^{k/8}$, this yields
\[
                    |\psi(x)-x|\le x2^{-k/8}.               \tag{A8}
\]
The $h/2$ estimate is valid because
$h/2\le x2^{-k/4}$ and $2^{\lfloor k/4\rfloor-k/2}\le2^{-k/4}$.
For integer $x$, $k$ is obtained by bounded search from its binary length:
choose the largest $j$ with $D_0 j^{10}\le\lfloor\log_2x\rfloor$.
Thus all rate functions displayed here are primitive recursive.

\paragraph{Status of the analytic estimate.}
The analytic estimates prove an explicit nonoptimal stretched-exponential
PNT error, equivalent to $O(x\exp(-c(\log x)^{1/10}))$ with a positive
explicit $c$. The classical square-root logarithmic rate needs a stronger
zero-free strip and is not proved by these notes. An iPRA formalization
still has to instantiate the indicated finite rational complex-arithmetic
identities and their error budgets; the fact that the numerical bounds are
primitive recursive is not, by itself, a proof that the argument is in iPRA.

\section{From the smoothed estimate to the stated arithmetic modulus}
\label{sec:theta-bridge}

Recall
\[
 \psi(x)=\sum_{p^j\le x}\log p,
 \qquad \vartheta(x)=\sum_{p\le x}\log p.
\]
Only powers with exponent at least two contribute to their difference.
For each such prime $p\le\sqrt x$, the total contribution is at most
$\log x$. There are at most $\lfloor\sqrt x\rfloor$ possible primes.
Consequently, by a finite sum,
\begin{equation}\label{eq:prime-powers-gap}
 0\le\psi(x)-\vartheta(x)\le\sqrt x\log x.
\end{equation}
The upper bound does not need an estimate for the distribution of primes.

Put $R=\log x$. For $R\ge16$, $\log R\le R/4$, so
$R e^{-R/2}\le e^{-R/4}$. On the dyadic scale of the preceding
section, $R\ge2^{98}k^{10}$, and therefore
\[
 \frac{|\psi(x)-\vartheta(x)|}{x}
 \le R e^{-R/2}\le e^{-R/4}\le2^{-k/8}.
\]
Combining this with (A8) gives
\begin{equation}\label{eq:theta-dyadic}
 |\vartheta(x)-x|\le2x\,2^{-k/8}.
\end{equation}
If a rational expression is desired at every intermediate step, replace
$2^{-k/8}$ by $2^{-\lfloor k/8\rfloor}$; the following constants have
much more than the one extra bit needed for this replacement.

\begin{theorem}[A simple primitive recursive modulus for $\vartheta$]
\label{thm:theta-modulus-established}
For $C=2^{240}$ and every natural $r$,
\begin{equation}\label{eq:theta-established}
 n\ge2^{C(r+1)^{10}}
 \quad\Longrightarrow\quad
 |\vartheta(n)-n|\le2^{-r}n.
\end{equation}
Here and below a real-valued inequality means its rational enclosure
version at arbitrary prescribed accuracy, with the finite interpretation
of Section~\ref{sec:foundations}.
\end{theorem}

\begin{proof}
Let $\ell=\lfloor\log_2n\rfloor$ be obtained from the binary length,
and let $k$ be the largest integer with $2^{100}k^{10}\le\ell$.
Bounded search up to $\ell$ supplies it. The hypothesis gives
\[
 k\ge2^{14}(r+1)\ge16384>8192,
\]
since $2^{100}(2^{14}(r+1))^{10}=2^{240}(r+1)^{10}$.
Maximality gives $\ell<2^{100}(k+1)^{10}$, and the latter is an
integer, so $\log_2n<\ell+1\le2^{100}(k+1)^{10}$.
Thus (A8) and \eqref{eq:theta-dyadic} apply. Finally
$2\cdot2^{-\lfloor k/8\rfloor}\le2^{-r}$, proving the assertion.
All comparisons selecting $\ell,k$ are integer comparisons.
\end{proof}

The next section proves an abstract transfer theorem, then applies it
with the modulus just established. Its prime-counting hypothesis is
therefore discharged here; it is not an extra assumption in the main
result. In particular it will give, for $D=2C$,
\[
 n\ge2^{D(m+5)^{10}}\ \Longrightarrow\
 |\pi(n)-\Li_2(n)|\le2^{-m}\frac n{\log n}.
\]
Because $b(n)\le2\log_2n\le4\log n$ for $n\ge2$, its instance at
$m+2$, followed by \eqref{eq:Li-error}, proves
Theorem~\ref{thm:object}. The transfer proof actually retains a factor
$17/32$ on its right-hand side, so approximating its analytic names
cannot consume the entire claimed bound.

% Finite transfer theorems; included by the main manuscript.
% This file deliberately has no preamble, and can be \input into the main report.
% All prime-number information is isolated in hypothesis (H_theta).

\section{A finite transfer theorem for the logarithmic integral}
\label{sec:finite-transfer}

Put
\[
 \vartheta(n)=\sum_{\substack{p\le n\\p\text{ prime}}}\log p,
 \qquad \pi(n)=\#\{p\le n:p\text{ prime}\},
 \qquad \operatorname{Li}_2(x)=\int_2^x\frac{dt}{\log t}.
\]
The lower limit $2$ removes the singularity at $1$.  Changing to the
principal-value convention $\operatorname{li}(x)$ changes the answer by a
fixed computable constant; it changes none of the error classes below.
All statements about $n$ in this section concern integers $n\ge2$.
The integrals and logarithms are abbreviations for their explicit rational
approximations, as described in the foundational section.

\subsection{The exact finite identity}

For any finite sequence $a_2,\ldots,a_n$, let
$A_k=\sum_{j=2}^k a_j$.  For any $b_2,\ldots,b_n$, distributivity and
telescoping give
\begin{equation}\label{eq:finite-abel}
 \sum_{j=2}^n a_jb_j
   =A_nb_n+\sum_{k=2}^{n-1}A_k(b_k-b_{k+1}).
\end{equation}
Indeed, the coefficient of $a_j$ on the right is
$b_n+\sum_{k=j}^{n-1}(b_k-b_{k+1})=b_j$.
This is an identity of finite sums over an arbitrary commutative ring;
its proof uses only induction on the length of the sum.

Set $a_j=\log j$ when $j$ is prime and $a_j=0$ otherwise, and set
$b_j=1/\log j$.  Write
\[
 E_k=\vartheta(k)-k,\qquad
 \Delta_k=\frac1{\log k}-\frac1{\log(k+1)},\qquad
 J_n=\frac n{\log n}+\sum_{k=2}^{n-1}k\Delta_k.
\]
Then \eqref{eq:finite-abel} is exactly
\begin{equation}\label{eq:pi-discrete-transfer}
 \pi(n)-J_n=\frac{E_n}{\log n}
                  +\sum_{k=2}^{n-1}E_k\Delta_k.
\end{equation}
Telescoping the definition of $J_n$ also gives
\begin{equation}\label{eq:discrete-reference}
 J_n=\frac2{\log2}+\sum_{k=3}^n\frac1{\log k}.
\end{equation}
Since $t\mapsto1/\log t$ is decreasing on $[2,n]$, the right and left
Riemann sums on the unit partition imply
\[
 \sum_{k=3}^n\frac1{\log k}
 \le \operatorname{Li}_2(n)
 \le \sum_{k=2}^{n-1}\frac1{\log k}.
\]
Consequently,
\begin{equation}\label{eq:discrete-Li-error}
 0\le J_n-\operatorname{Li}_2(n)\le\frac2{\log2}\le4.
\end{equation}
The estimates $1/2\le\log2\le3/4$ suffice everywhere below.  They follow,
for example, by integrating $1/t$ on $[1,2]$, with the upper bound supplied
by the trapezoid inequality for the convex function $1/t$.

Two further bounds will be used:
\begin{equation}\label{eq:delta-and-trivial}
 0\le\Delta_k\le\frac1{k(\log k)^2},
 \qquad |E_k|\le3k\log k.
\end{equation}
For the first, use $\log(1+1/k)\le1/k$ and
$\log(k+1)\ge\log k$.
For the second, $0\le\vartheta(k)\le k\log k$ and
$k\le2k\log k$ imply
$|\vartheta(k)-k|\le\vartheta(k)+k\le3k\log k$.
No theorem about the density of primes is used here.

\subsection{Transfer of an arbitrary explicit modulus}

Suppose $T:\mathbb N\to\mathbb N$ is a primitive recursive function,
$T(r)\ge2$, and the following prime-number estimate has been established:
\begin{equation}\tag{$H_\vartheta$}\label{hyp:theta-modulus}
 k\ge T(r)\quad\Longrightarrow\quad
 |\vartheta(k)-k|\le2^{-r}k
 \qquad(r\in\mathbb N).
\end{equation}
This is a hypothesis of the abstract transfer lemma. In the application here
it is supplied by Theorem~\ref{thm:theta-modulus-established}; it is not an
extra assumption of the final theorem.

\paragraph{Theorem (explicit logarithmic-integral transfer).}
Define
\[
 T_\pi(m)=\max\left\{T(m+4)^2,\ 2^{16(m+5)}\right\}.
\]
Then
\begin{equation}\label{eq:pi-modulus}
 n\ge T_\pi(m)\quad\Longrightarrow\quad
 |\pi(n)-\operatorname{Li}_2(n)|
       \le 2^{-m}\frac n{\log n}.
\end{equation}

\paragraph{Proof.}
Let $L=\log n$, $\delta=2^{-(m+4)}$, and
$q=\lceil\sqrt n\rceil$.  The integer $q$ is obtained by bounded search
for the least $q\le n$ satisfying $n\le q^2$.
The hypothesis on $n$ gives $q\ge T(m+4)$.
For $2\le k<q$, \eqref{eq:delta-and-trivial} gives
$|E_k|\Delta_k\le3/\log k\le6$.
There are at most $\sqrt n$ such indices, so their total is at most
$6\sqrt n$.
For $q\le k<n$, hypothesis \eqref{hyp:theta-modulus} and
$\log k\ge L/2$ give
\[
 |E_k|\Delta_k\le\frac\delta{(\log k)^2}
                         \le\frac{4\delta}{L^2}.
\]
There are fewer than $n$ such indices.  Applying
\eqref{eq:pi-discrete-transfer} and \eqref{eq:discrete-Li-error},
\begin{equation}\label{eq:raw-transfer}
 |\pi(n)-\operatorname{Li}_2(n)|
 \le\delta\frac nL+4\delta\frac n{L^2}+6\sqrt n+4.
\end{equation}
We now check the numerical margins explicitly.  The condition
$n\ge2^{16(m+5)}$ implies $L\ge8(m+5)\ge40$.
For $L\ge16$, $\log L\le\sqrt L\le L/4$; hence
\[
 \frac{\sqrt n}{L}
   =\exp(L/2-\log L)
   \ge\exp(L/4)
   \ge\exp(2(m+5))\ge2^{m+5}.
\]
Also $4\le\sqrt n$ and $L\ge1$.  Thus the right side of
\eqref{eq:raw-transfer} is at most
\[
 \left(\frac5{16}+\frac7{32}\right)2^{-m}\frac nL
 =\frac{17}{32}\,2^{-m}\frac nL
 \le2^{-m}\frac nL.
\]
The unused margin $15/32$ is helpful when replacing the analytic notation
by rational enclosures.  This proves the theorem.

\paragraph{Why the elementary logarithmic inequality is legitimate.}
For completeness, $\log L\le\sqrt L$ for $L\ge1$ follows on writing
$v=\sqrt L$: the function $v-2\log v$ decreases on $[1,2]$ and increases
on $[2,\infty)$, and its value at $2$ is positive because $\log2<1$.
Equivalently, the same assertion follows from the corresponding finite
integral bounds; no compactness or choice is involved.

\subsection{A useful stretched-exponential specialization}

Suppose, for fixed positive integers $C,d$, that
\begin{equation}\label{eq:theta-power-modulus}
 n\ge2^{C(r+1)^d}\quad\Longrightarrow\quad
 |\vartheta(n)-n|\le2^{-r}n.
\end{equation}
Take $D=\max(2C,16)$.  The transfer theorem immediately gives
\begin{equation}\label{eq:pi-power-modulus}
 n\ge2^{D(m+5)^d}\quad\Longrightarrow\quad
 |\pi(n)-\operatorname{Li}_2(n)|\le2^{-m}\frac n{\log n}.
\end{equation}
For $d\ge3$ one may take the smaller $D=2C$, since
$2C(m+5)^d\ge16(m+5)$.  In particular this simplification applies to
$d=10$.

The usual analytic notation for the same conclusion is
\begin{equation}\label{eq:pi-stretched}
 |\pi(n)-\operatorname{Li}_2(n)|
 \le64\,\frac n{\log n}
       \exp\left(-\frac{(\log n)^{1/d}}{2D}\right)
 \qquad(n\ge2^{D6^d}).
\end{equation}
Here is a constructive derivation that does not select an integer by an
undecidable comparison of real numbers.  Let $\ell$ be the unique integer
with $2^\ell\le n<2^{\ell+1}$, and let $q$ be the largest integer with
$Dq^d\le\ell$.  Both are obtained by bounded integer search.  Then
$q\ge6$, and \eqref{eq:pi-power-modulus} applies with $m=q-5$.
Moreover, $\ell\ge6$ and $\log2\le3/4$ give
\[
 \log n<(\ell+1)\log2\le\tfrac34(\ell+1)\le\ell.
\]
Thus
\[
 q> (\ell/D)^{1/d}-1
      \ge(\log n)^{1/d}/D-1,
\]
where $D^{1/d}\le D$ was used.  Therefore
\[
 2^{-(q-5)}\le64\exp\bigl(-(\log2)(\log n)^{1/d}/D\bigr)
       \le64\exp\bigl(-(\log n)^{1/d}/(2D)\bigr),
\]
which proves \eqref{eq:pi-stretched}.
For the intended $d=10$ and $D=2C$, the exponential coefficient is
$1/(4C)$.
For the specific constant $D=2^{241}$ we can retain the tenth root
instead of replacing it by $D$. Since $D^{1/10}<2^{25}$ and
$\log2\ge1/2$, the same calculation gives the stronger explicit form
\begin{equation}\label{eq:fixed-stretched}
 |\pi(n)-\operatorname{Li}_2(n)|
 \le64\frac n{\log n}
       \exp\!\left(-2^{-26}(\log n)^{1/10}\right),
 \qquad n\ge2^{D6^{10}}.
\end{equation}
All thresholds derived below from the weaker generic coefficient remain
valid. This modest simplification changes no arithmetic construction.


\subsection{Every inverse-logarithmic order, with a primitive recursive threshold}

The stretched-exponential estimate implies much more than
$\pi(n)\sim\operatorname{Li}_2(n)$.  For $k\in\mathbb N$ set
\[
 M(k)=1+16D^2d^2k^2+24D,
 \qquad N(k)=2^{2M(k)^d}.
\]
Then
\begin{equation}\label{eq:all-log-orders}
 n\ge N(k)\quad\Longrightarrow\quad
 |\pi(n)-\operatorname{Li}_2(n)|
        \le\frac n{(\log n)^{k+1}}.
\end{equation}
To verify this, put $L=\log n$ and $u=L^{1/d}$.
The lower bound on $n$ gives $u\ge M(k)$.  In particular,
$u\ge24D$ and $u\ge16D^2d^2k^2$.
Since $\log u\le\sqrt u$ for $u\ge1$,
\[
 \log64+kd\log u
 \le6+kd\sqrt u
 \le\frac{u}{4D}+\frac{u}{4D}
 =\frac{u}{2D}.
\]
Exponentiating yields
$64\exp(-u/(2D))\le u^{-kd}=L^{-k}$.
The same lower bound on $n$ is stronger than $n\ge2^{D6^d}$, because
$M(k)\ge24D$.  Combining with \eqref{eq:pi-stretched} proves
\eqref{eq:all-log-orders}.
When $d=10$ and $D=2C$, this is the explicit formula
\[
 M(k)=1+6400C^2k^2+48C.
\]
The thresholds are intentionally coarse.  Their purpose is to be
transparent, uniform primitive recursive witnesses.

\subsection{The full factorial asymptotic expansion}

For an integer $r\ge1$ define
\[
 S_r(n)=n\sum_{j=0}^{r-1}\frac{j!}{(\log n)^{j+1}},
 \qquad B_r=1+r!\bigl(1+2^{r+2}\bigr).
\]
Repeated integration by parts, with every boundary term retained, gives
\begin{equation}\label{eq:Li-expansion-exact}
 \operatorname{Li}_2(n)
  =S_r(n)
   -2\sum_{j=0}^{r-1}\frac{j!}{(\log2)^{j+1}}
   +r!\int_2^n\frac{dt}{(\log t)^{r+1}}.
\end{equation}
For clarity, the one-step identity behind this formula is
\[
 \int_2^n\frac{dt}{(\log t)^s}
 =\frac n{(\log n)^s}-\frac2{(\log2)^s}
       +s\int_2^n\frac{dt}{(\log t)^{s+1}}.
\]
It follows by integrating the explicit derivative of
$t/(\log t)^s$; induction on $s$ yields
\eqref{eq:Li-expansion-exact}.

Suppose $L=\log n\ge16(r+2)^2$.  Splitting the remaining integral at
$\sqrt n$ gives
\[
 \int_2^n\frac{dt}{(\log t)^{r+1}}
 \le 2^{r+1}\sqrt n+\frac{2^{r+1}n}{L^{r+1}}
 \le\bigl(1+2^{r+1}\bigr)\frac n{L^{r+1}}.
\]
For the second inequality it is enough that
$(r+1)(\log2+\log L)\le L/2$.
Writing $v=\sqrt L\ge4(r+2)$, the left side is at most
$(r+1)(1+v)\le(v/4)(2v)=L/2$.
The boundary constant satisfies
\[
 2\sum_{j=0}^{r-1}\frac{j!}{(\log2)^{j+1}}
 \le2\sum_{j=0}^{r-1}j!2^{j+1}
 \le2^{r+1}r!.
\]
Also $n/L^{r+1}\ge1$ under the same lower bound on $L$.
It follows that
\begin{equation}\label{eq:Li-expansion-bound}
 |\operatorname{Li}_2(n)-S_r(n)|
 \le r!\bigl(1+2^{r+2}\bigr)\frac n{L^{r+1}}.
\end{equation}
Combining with \eqref{eq:all-log-orders}, we obtain the explicit theorem
\begin{equation}\label{eq:pi-expansion-bound}
 \begin{split}
 n&\ge2^{\max\{2M(r)^d,\,32(r+2)^2\}}\\
 &\Longrightarrow\quad
 \left|\pi(n)-n\sum_{j=0}^{r-1}\frac{j!}{(\log n)^{j+1}}\right|
 \le B_r\frac n{(\log n)^{r+1}}.
 \end{split}
\end{equation}
Thus every fixed number of terms of the usual factorial expansion is
valid, and both the coefficient and the threshold are given by primitive
recursive expressions in the requested order.

There are also effective one-sided statements.  Apply
\eqref{eq:pi-expansion-bound} with $r+1$ terms and additionally require
$\log n\ge2B_{r+1}/r!$.  Since
$S_{r+1}(n)-S_r(n)=r!n/(\log n)^{r+1}$, this gives
\[
 \frac{r!}{2}\frac n{(\log n)^{r+1}}
 \le\pi(n)-S_r(n)
 \le\frac{3r!}{2}\frac n{(\log n)^{r+1}}.
\]
For example, the extra natural-number condition
$n\ge2^{4B_{r+1}}$ is sufficient.  Consequently every fixed truncation
$S_r$ eventually underestimates $\pi$, and its first omitted term gives
the leading size of the error.  This assertion is about fixed
truncation order; it does not assert convergence of the factorial series.

\subsection{A constructive prime-in-an-interval consequence}

Under \eqref{eq:pi-power-modulus}, let $m\in\mathbb N$ and
$n\ge2^{D(m+9)^d}$.  Put $h=\lceil n/2^m\rceil$ and
$\delta=2^{-(m+4)}$.  Both $n$ and $n+h$ satisfy the estimate with
precision $m+4$, and $1\le h\le n$.
The logarithmic-integral increment satisfies
\[
 \operatorname{Li}_2(n+h)-\operatorname{Li}_2(n)
 \ge\frac h{\log(n+h)}\ge\frac h{2\log n}.
\]
The sum of the two endpoint error bounds is at most
$3\delta n/\log n$, because $n+h\le2n$ and
$\log(n+h)\ge\log n$.
As $h\ge n/2^m=16\delta n$, we conclude
\[
 \pi(n+h)-\pi(n)
 \ge\frac{h/2-3\delta n}{\log n}
 \ge\frac{5\delta n}{\log n}>0.
\]
Therefore a prime exists in $(n,n+\lceil n/2^m\rceil]$.
It is found by bounded search in that explicitly given interval.
This is a witness-producing consequence, not merely a negated
statement about empty intervals.  There is no hidden search for a
precision at which strict positivity becomes visible: since
$\log n\le n$, the displayed lower bound is at least $5\delta$.
A predetermined precision larger than $m+10$ therefore suffices to
certify positivity with rational enclosures.

\subsection{What the transfer does and does not prove about iPRA}

The prime-number hypothesis must itself be proved in iPRA; ordinary
computability of $T$ is not a substitute for such a proof.
The present section supplies the following concrete reduction.

\begin{enumerate}
 \item Equation \eqref{eq:finite-abel} is proved for coded finite rational
 lists by quantifier-free induction on list length.  It remains exact if
 the $\log j$ are replaced by positive rational approximations and each
 $b_j$ is defined as the reciprocal of that same approximation.
 \item Primality, $\pi(n)$, finite sums, integer square roots, binary
 length, fixed integer powers, and the thresholds displayed above are
 primitive recursive.  Bounded search is always given an explicit
 natural-number upper bound.
 \item The bounds for logarithms follow from finite rational sums and
 geometric-series tails.  Riemann-sum inequalities for $\operatorname{Li}_2$
 are finite inequalities with an explicit mesh error.  On $[2,n]$,
 $1/\log t\le2$ and $|\frac{d}{dt}(1/\log t)|\le2$ suffice for a uniform
 rational integration modulus.
 \item Each displayed real inequality can be replaced by the assertion
 that every precision-$s$ rational enclosure satisfies the desired
 inequality with an additional explicit $2^{-s}$ allowance.
 For a given $n,m,s$, all sums and all enclosures involved have
 primitive recursive bounds.  One first proves the corresponding
 bounded rational inequality by quantifier-free induction, and only
 then universally generalizes $n,m,s$.  One does not apply induction
 to a formula containing a quantifier over all approximation precisions.
 \item The integration-by-parts identity for the factorial expansion is
 established by the finite telescoping identity for the explicit
 derivatives, together with a Taylor-remainder bound for each rational
 mesh interval.  The derivative bounds involve only rational powers
 of a positive lower bound for $\log t$; the mesh can therefore be
 chosen by a primitive recursive formula.  Induction in the expansion
 order is again performed on the finite rational identity with its
 explicit error allowance.
\end{enumerate}

For a particularly simple all-arithmetic conclusion, let $b(n)$ be the
binary length of $n$, and let $A(n,s)\in\mathbb Q$ be an explicitly
constructed approximation to $\operatorname{Li}_2(n)$ with
$|A(n,s)-\operatorname{Li}_2(n)|\le2^{-s}$.
Since $b(n)\le4\log n$, \eqref{eq:pi-modulus} at precision $m+2$ implies
\[
 n\ge T_\pi(m+2)\quad\Longrightarrow\quad
 |\pi(n)-A(n,s)|
       \le\frac{n}{2^m b(n)}+2^{-s}.
\]
After clearing positive integer denominators this is a decidable
arithmetical inequality.  This is the recommended final object-language
statement for the Lean-to-iPRA development.

\paragraph{Lean dependency outline.}
First implement finite Abel summation over coded rational lists and
prove its coefficient-telescoping identity.  Next prove the prime
specialization, the discrete-reference formula, and the three bounds
\eqref{eq:discrete-Li-error}--\eqref{eq:delta-and-trivial}.  The proof of
\eqref{eq:raw-transfer} then consists only of splitting a finite sum at
the bounded integer square root.  The remaining modulus calculations
are elementary natural-number and rational inequalities.  The logical
separation should be reflected in Lean: the input theorem about
$\vartheta$ is an explicit hypothesis or an imported proved theorem;
it must not be hidden behind an axiom named as though it were already
established.  Proving the corresponding theorem using Mathlib's real
numbers is a useful reference development, but a separate derivation
or checked translation into the object theory is required before
claiming an iPRA theorem.

% Independently derived finite-arithmetic appendix, prepared by repo_foundations.
% Intended as an input section, not a standalone document.
\section{Finite contour cancellation and uniform arithmetic schedules}
\label{sec:finite-contour-audit}

This section records concrete replacements for the analytic notation used
in the contour argument. In particular, it does not infer provability from
the mere existence of a primitive recursive approximation algorithm.

\subsection{A strictly finite Cauchy lemma}

Use the taxicab norm $\lVert u+iv\rVert_1=|u|+|v|$ on rational complex
numbers. Let a rational rectangle have width $W$ and height $H$, and fix
an integer $R\ge W+H$. Suppose that, on the rectangle, a function $G$
has rational approximation procedures for $G$ and $G'$ with error at most
$\eta$, and a proved uniform Taylor inequality
\[
 \lVert G(z)-G(a)-G'(a)(z-a)\rVert_1
       \le M\lVert z-a\rVert_1^2
\]
whenever the segment from $a$ to $z$ stays in the rectangle. The displayed
inequality means its finite rational enclosure version with the indicated
evaluation errors. Here $M\ge1$ is an integer bound, obtained from the
concrete formula for $G$; it is not selected by compactness.

Divide each side into $N$ equal parts. Each of the $N^2$ small rectangles
has taxicab diameter $d=(W+H)/N$ and perimeter $2d$. On each cell choose
its lower left vertex $a$ and the rational affine function
\[
 A_a(z)=G_\eta(a)+G'_\eta(a)(z-a).
\]
For every cell vertex $v$ the rational residual satisfies
\[
 \lVert G_\eta(v)-A_a(v)\rVert_1
       \le Md^2+\eta(2+d).
\]
The trapezoidal boundary sum of an affine complex function is exactly
zero: the constant part telescopes, while the linear part uses
$(u+v)(v-u)=v^2-u^2$. Hence the cell boundary sum is at most
$2d[Md^2+\eta(2+d)]$. Every internal edge is used twice with opposite
orientations and the same rational vertex values. It cancels exactly.
The remaining outer boundary trapezoidal sum $C_{N,\eta}$ therefore obeys
\begin{equation}\label{eq:finite-rectangle}
 \lVert C_{N,\eta}\rVert_1
 \le \frac{2MR^3}{N}+4RN\eta+2R^2\eta.
\end{equation}
This proves the finite lemma using only distributivity, the triangle
inequality, finite summation, and exact incidence cancellation. In
particular the evaluation allowance includes the growing internal
perimeter; omitting its factor $N$ would be incorrect.

For requested error $2^{-q}$ take
\[
 N=2^{q+6}(M+1)(R+1)^3,
 \qquad
 \eta\le\frac{2^{-q-8}}{(R+1)^2(N+1)}.
\]
The three terms of \eqref{eq:finite-rectangle} are respectively at most
$2^{-q-5},2^{-q-6},2^{-q-7}$, so their sum is smaller than $2^{-q-4}$.
If $b(u)$ is the least integer with $u\le2^{b(u)}$, a sufficient sample
precision is
\[
 q_{\rm ev}=q+8+2b(R+1)+b(N+1).
\]
The search for $b(u)$ is bounded by $u$. Enlarge $M$, if needed, to also bound $|G'|$. This gives the usual
finite rectangle-refinement estimate for the contour integral name; its
mesh error is at most a constant times $MR^2/N$. Increasing $q$ by four
at the start absorbs this additional comparison. Thus the integral
notation abbreviates a theorem about finite rational tables at every
requested precision, rather than an appeal to a general integration
theorem.

\subsection{A concrete derivative bound after pole subtraction}

For the PNT contour put
\[
 A(s)=\frac{y^{s+1}}{s(s+1)},\qquad
 G(s)=-\frac{\zeta'(s)}{\zeta(s)}A(s)-\frac{y^2}{2(s-1)}.
\]
The removable value at $s=1$ is evaluated by its regular expression below.
One deliberately generous usable bound, in the taxicab norm, is
\[
 M=2^{512}(\lceil y\rceil+1)^8(T+4)^{40}.
\]
It bounds the magnitudes of $G,G',G''$ and therefore its Taylor remainder,
with room for fixed norm-conversion factors. Here is a direct check.

Outside $|s-1|<1/128$, Euler summation with lower cutoff $1$ gives,
for $0\le j\le3$, bounds for $\zeta^{(j)}$ using
\[
 \left|\int_1^\infty\{u\}u^{-s-1}(\log u)^j\,du\right|
       \le\frac{j!}{\sigma^{j+1}}\le j!2^{j+1}.
\]
The pole derivatives are at most $j!128^{j+1}$. Thus each
$|\zeta^{(j)}|$ is bounded by $Z=2^{40}(T+3)$.
The proved strip apartness gives
$|1/\zeta|\le J=2^{71}(T+3)^7$. On the remaining low-ordinate band the same reciprocal bound follows from
$|1/\zeta|=|s-1|/|w|\le(1/32)/(7/8)<1\le J$.
With $U=ZJ$, ordinary quotient
identities bound $|F|,|F'|,|F''|$, for $F=-\zeta'/\zeta$, respectively
by $U,2U^2,6U^3$. The identities
\[
 \frac1{s(s+1)}=\frac1s-\frac1{s+1},\qquad
 \frac{d^j}{ds^j}y^{s+1}=(\log y)^j y^{s+1}
\]
show that $|A^{(j)}|\le2^6 y^6$ for $0\le j\le3$ on
$1/2\le\sigma\le2$, since $\log y\le y$ and $y\ge2$.
Expanding $G''$ by the product rule gives a bound smaller than
$2^{360}(T+3)^{24}y^6$.

Near the pole write $w=(s-1)\zeta=1+(s-1)h$ and
$g=-w'/w$. On the rational band
$|\sigma-1|\le1/64$, $|t|\le1/64$, the previously proved estimates
give $|w|\ge7/8$. Euler's formula bounds the first three derivatives
of $h$ by $5,10,40,240$, respectively. Hence the first three
derivatives of $w$ are bounded by $6,22,128$, and
$|g|,|g'|,|g''|\le2^{13}$. On this band use the stable expression
\[
 G(s)=\int_0^1 A'(1+u(s-1))\,du+g(s)A(s).
\]
Its second derivative is bounded directly by the displayed bounds for
$A^{(j)}$ and $g^{(j)}$; it is much smaller than the proposed $M$.
The contour can be partitioned at the rational lines $t=\pm1/64$.
Thus choosing between the two formulas requires only rational comparisons,
not a decision about an unknown complex zero. The tiny horizontal
distance $\delta_T$ from the pole causes no numerical instability in
the regular expression.

The simple-pole integral is also finite. Excise a rational square centred
at $1$ of half-side $\rho=\tfrac12\min\{\delta_T,c-1,1/64\}$,
which is positive and fits inside both vertical sides; cancel the rational grid between it
and the outer rectangle, and compute its four sides. Their real parts
cancel in pairs and their imaginary parts add to
$8i\int_0^1du/(1+u^2)=2\pi i$, where this identity defines the same
explicit rational integral name of $\pi$ as usual. In particular
$2\le\pi\le4$, directly from the integrand bounds. This proves the
residue calculation without an additional infinitary residue theorem.

\subsection{Fixed rational logarithm and exponential algorithms}

For a rational $u\in[1,2]$, put $z=(u-1)/(u+1)\in[0,1/3]$.
The finite sum
\[
 \ell_N(u)=2\sum_{j=0}^{N-1}\frac{z^{2j+1}}{2j+1}
\]
has tail at most $9^{-N}$, since
$2\sum_{j\ge N}|z|^{2j+1}\le(3/4)9^{-N}$.
For any positive rational $a/b\ge1$, bounded integer search finds
$j$ with $2^j\le a/b<2^{j+1}$. Write $a/b=2^ju$ and approximate
its logarithm by $j\ell_N(2)+\ell_N(u)$. The error is at most
$(j+1)9^{-N}$. For example $N=p+j+2$ gives error smaller than
$2^{-p}$; for $a/b<1$ use its reciprocal and change sign.
The identities for these names follow from finite geometric-series
identities or rational rectangle refinements, with these same tails.

For a rational complex $z$ with $\lVert z\rVert_1\le B$, where $B\ge1$ is an integer,
use the Taylor polynomial $\sum_{j<N}z^j/j!$. At index $2B$ its
absolute coefficient is at most $B^B\le2^{B^2}$, since the last $B$
factors of $(2B)!$ are at least $B$. Every subsequent ratio is at most
$1/2$. Thus
\[
 \sum_{j\ge N}\frac{B^j}{j!}
       \le2^{1+B^2+2B-N}\qquad(N\ge2B).
\]
Taking $N=B^2+2B+p+2$ gives error at most $2^{-p-1}$.
The case $B=0$ is exact. The derivative and multiplication identities
are obtained by differentiating the finite polynomial and applying the
same explicit tails. When the input itself is approximate, the series
gives the Lipschitz bound $4^{B+1}$ on the ball of radius $B$; the
input precision is chosen first to absorb this displayed factor.

For Euler summation and its first three derivatives, truncation at
$V=2^{2v}$ has a bound of the form
\[
 2^{12}(T+3)(v+1)^3 2^{-v}.
\]
Repeated integration by parts gives this explicitly from
\[
 \int_V^\infty u^{-\sigma-1}(\log u)^jdu
 =V^{-\sigma}\sum_{h=0}^j
    \frac{j!}{(j-h)!}\frac{(\log V)^{j-h}}{\sigma^{h+1}}.
\]
For instance $v=4(p+T+32)^2$ is sufficient for error $2^{-p}$.
The inequality $(v+1)^3\le2^{v/2}$ used here holds for $v\ge64$
and has an elementary finite-induction proof. Split the remaining
integral at all integers up to $V$. On each interval its fractional-part
factor is affine, and a crude derivative bound is
$K=2^{16}(T+3)(V+1)^4$. A rational mesh of step $1/Q$ with
$Q\ge2^{p+4}K(V+1)$ therefore has total quadrature error at most
$2^{-p-4}$. Evaluate its exponential and logarithm terms with enough
additional precision to absorb its finite coefficient sum. All cutoffs
are displayed bounded arithmetic constructions.

For a finite Euler product at real part $1+d$, $d>0$ rational, its
distance from the Dirichlet series is at most $P^{-d}/d$ when all
primes up to the integer $P$ are used. Choose
$P=2^V$ with $V\ge(p+d^{-1}+10)/d$. Then this tail is below $2^{-p}$.
Each factor has denominator apartness at least $1/2$. The real
logarithm of its squared modulus may be expanded with the elementary
geometric tail, and the identity
$3+4\cos v+\cos2v=2(1+\cos v)^2$ applies to every finite summand.
Choose each factor's input evaluation tolerance smaller than the desired
final error divided by $8P4^P[2+2/d]^{16}$. Reciprocal evaluation loses
at most a factor $8$, and a product of at most $P$ factors of magnitude
at most $2$ has error at most $P4^P$ times the factor error when all
factor errors are at most $1$. The remaining displayed multiplier
absorbs the eight-factor positivity comparison. This gives the finite version of the Euler-product
positivity inequality with explicit slack.

\subsection{Uniform expansion order and the induction discipline}

The factorial expansion need not be proved by induction on a statement
quantifying over real precisions. For the requested integer order $r$,
form the single finite expression
\[
 S_r(t)=t\sum_{j<r}\frac{j!}{(\log t)^{j+1}}.
\]
The derivative identity
\[
 S'_r(t)=\frac1{\log t}-\frac{r!}{(\log t)^{r+1}}
\]
is exact finite telescoping: the terms with adjacent indices cancel.
On $2\le t\le n$, an integer bound for $|S''_r|$ is
$B_r^*=2^{r+4}(r+2)^4r!$. Telescope over a rational mesh of step
at most $2^{-m}$, with $m=p+B_r^*+n+4$. Its total Taylor error is
at most $B_r^*n2^{-m}<2^{-p-3}$.
On $t\ge2$ clip each rational logarithm sample upward to $1/2$;
the true value is at least $1/2$, so clipping cannot increase its error.
Consequently every sampled reciprocal has the same positive lower
denominator bound used in the derivative estimates.
All logarithm powers and factorials are evaluated by fixed PR programs
with the parameter $r$; no universal evaluator for arbitrary PR codes
is being assumed. Logarithm evaluation tolerances can, for example,
be taken below $2^{-p-10}/[(n+1)(B_r^*+1)]$ to absorb the finite
coefficient derivatives. This yields the expansion identity uniformly
in $r$ after denominators are cleared.

In every application, fix the outer numerical parameters, the requested
precision, all cutoffs, and all rational denominator bounds before
induction. A loop invariant is an equality or inequality between
primitive recursive integer expressions, or the Boolean conjunction of
such checks over the already processed prefix. Quantifier-free induction
proves the row invariant and its finite fold. Universal introduction is
used only after that derivation. In particular, there is no induction on
``all precisions are correct'', ``the function is holomorphic'', or
``all sufficiently large arguments satisfy the estimate''. The analytic
notation is justified by the proved finite identities and their explicit
error budgets, rather than by a generic assertion that computability
implies arithmetic provability.

\section{Lean formalization: exact targets and a completion route}
\label{sec:formalization}

\subsection{Two different correctness questions}

There are two tasks. The first is to formalize the mathematical inequalities
using Lean's natural numbers, rationals, and, if convenient for a reference
development, Mathlib's real and complex numbers. The second is to construct
an object-language derivation in \ipra. The first task does not entail the
second. Classical tactics in a reference development must not be silently
exported to the object theory.

The repository provides the syntax \texttt{PR}, \texttt{PTm},
\texttt{Fm}, the intuitionistic proof relation \texttt{Prf}, and the
axioms \texttt{PRAAx}. The required final expression has the shape
\[
 \texttt{Prf PRAAx [] mainFormula},
\]
where \texttt{mainFormula} is the universal closure of
\eqref{eq:integer-matrix} under the threshold antecedent.
Only after constructing that proof may one apply
\texttt{pra\_iff\_native} to obtain its native-calculus counterpart.
The separate bridge file in the archive demonstrates this use with an
\emph{explicitly supplied proof argument}; it does not supply the PNT proof.

\subsection{The implementation sequence}

\begin{longtable}{@{}p{0.19\textwidth}p{0.73\textwidth}@{}}
\toprule
Stage & Concrete proof obligation and endpoint\\
\midrule
\endhead
Arithmetic coding & Define signed integers and positive-denominator rationals;
prove cross-multiplication order, absolute value, finite summation, powers,
factorial, binary length, and bounded search by quantifier-free induction.
Represent every function by a fixed \texttt{PR} expression.\\[3pt]
Primes & Define bounded divisibility and primality, the count $\Pi$,
prime-power detection, and rational approximations to $\psi,\vartheta$.
Verify bounded factorization by the invariant that the product already
extracted times the residual is the original integer. Prove the elementary
unique-factorization identities used in finite Euler products.\\[3pt]
Elementary names & Implement \eqref{eq:log-algorithm}, the finite exponential,
and \eqref{eq:Li-algorithm}. Prove their displayed rational enclosure bounds.
Use explicit positive lower bounds for every reciprocal.\\[3pt]
Finite calculus & Prove interval telescoping with Taylor remainders, the
rectangle identity \eqref{eq:finite-rectangle}, and the explicit square
integral for a simple pole. All mesh sizes and sample accuracies are inputs
or prescribed PR terms.\\[3pt]
Zeta construction & Implement finite Euler summation, with (E3) as its
error budget, and the finite Euler products. Prove (E5)--(E7) with their
cutoffs. This includes the logarithmic derivative identity, not an axiom
asserting it.\\[3pt]
Apartness & Prove the polynomial cosine identity and the finite product
bound $\ge1/2$. Derive the rational-parameter inequalities (A3r), and
the regular expression (A4) on the small-height band.\\[3pt]
Counting formula & Prove the scalar kernel (E8), then its finite-sum
specialization (A5). Apply finite contour cancellation after subtracting the
known pole and obtain (A6) with the stated constants.\\[3pt]
Dyadic estimate & Derive (A7), the two finite-difference inequalities,
(A8), and Theorem~\ref{thm:theta-modulus-established}. Integer roots and
floor divisions are implemented by bounded search.\\[3pt]
Transfer & Instantiate finite Abel summation, use the split at the bounded
integer square root, and prove the factor $17/32$ estimate. Combine with
the error of $\mathcal A$ and clear denominators.\\[3pt]
Export & Universally close the quantifier-free matrix, obtain the actual
\texttt{Prf PRAAx [] mainFormula}, then use the repository's proved
native translation. Audit the dependency tree for added object axioms.\\
\bottomrule
\end{longtable}

This order is a dependency order. It does not require formalizing the
functional equation of $\zeta$, Hadamard factorization, an argument principle
for arbitrary holomorphic functions, an enumeration of zeros, or a general
Tauberian theorem. The main estimate uses one explicitly given pole and
one finite rectangle identity. Existing work on PNT in other proof
assistants or classical Lean blueprints can guide a reference proof, but
their proof-theoretic strength is not automatically that of \ipra.

\subsection{A specific error-budget convention}

It is useful to make the intended compilation mechanically checkable.
An arithmetic lemma carries the precision request and a rational error
allowance as explicit parameters. When a derived expression has a known
Lipschitz multiplier $K$ and uses at most $J$ input samples, evaluate
each sample with error at most
\[
 \frac{2^{-v}}{16(J+1)(K+1)}.
\]
The required number of bits is
$v+4+b(J+1)+b(K+1)$, increased by two if the input interval itself is
rounded. For products, $K$ is obtained from the known magnitude bounds
and the product rule. For the contour lemma, use its more precise
allowance, which accounts for the internal perimeter.

For the last inference of Theorem~\ref{thm:object}, the transfer proof
leaves $15/32$ of the prime-counting allowance, and the Li algorithm
uses less than $19/64$ of the allowance $2^{-v}$. Allocate at most
$2^{-v-4}$ to the auxiliary analytic computations after their displayed
amplification factors. Every cutoff is then a concrete PR expression in
$n,m,v$. The finite invariants establish the remaining integer inequality;
they are not tests of finitely many small instances of that inequality.

\subsection{What the included Lean files establish}

The included source files give precise target definitions, small finite
algebraic proof candidates, and the existing-library bridge. They use no
new axiom called ``PNT'' and do not conceal the analytic input in an
assumption presented as a theorem. They have \emph{not been compiled in
this workspace}: a Lean toolchain was not available. Consequently no
kernel-checking claim is made for these new files. The pinned repository
declares Lean and Mathlib version \texttt{v4.32.0}; the checked build of
that repository was not rerun here.

The completion criterion is stronger than ``the files contain no
\texttt{sorry}''. A finished development must define the PR codes for
the arithmetic programs, identify \texttt{mainFormula} with the actual
matrix, construct its \texttt{Prf PRAAx} derivation, and audit all
dependencies. The present archive supplies a developed mathematical proof
and a concrete route to those objects, as requested, rather than asserting
that this work is already mechanized.

\section{Comparison with the linked strengthenings and inspected drafts}
\label{sec:scope}

The logarithmic integral has two common normalizations:
$\Li_2(x)=\int_2^xdt/\log t$ and the principal-value $\operatorname{li}(x)$.
They differ by $\operatorname{li}(2)$, a fixed computable real. The explicit
arithmetic target in this note uses only $\Li_2$. All asymptotic error
classes asserted here are unchanged by a fixed additive constant;
their numerical thresholds need enlargement if a different normalization
is substituted. We do not transfer a small-argument numerical constant
between the two normalizations without accounting for this difference.

The standard $\log^{-9}T$ argument used here is an established analytic
method; a corresponding classical proof is discussed in the
Kontorovich project blueprint \cite{Kontorovich}. The contribution of this
note is its explicit rational apartness bounds, coarse integer moduli,
finite error accounting, and the specified weak-arithmetic target. The
classical square-root-logarithm estimate requires a stronger zero-free
region; the Vinogradov--Korobov error requires stronger analytic estimates
again. Neither stronger input is proved in this package. The bounds
$O(\sqrt x\log x)$ and Schoenfeld's sharp conditional inequality also
must retain their Riemann-hypothesis assumption. The introductory scope
table is therefore part of the statement of the result, not an omitted
technical qualification. See \cite{Tao,Trudgian} for stronger classical
estimates and their analytic dependencies.

The inspected archive
\texttt{docs/drafts/pnt\_ipra\_constructive\_2026-11-14.zip}
claims a modulus of the shape
\[
 \left|\frac{\psi(x)}x-1\right|\le\frac1q
 \quad\text{for }x\ge2^{\,2^{(2^{20}q^3)^2+1}}.
\]
That stated modulus alone does not yield an error smaller than every
inverse power of $\log x$. The new proof therefore does not assume that
draft as its analytic input. Its Lean artifacts, and the inspected Cauchy
and zeta drafts, explicitly describe themselves as uncompiled outlines.
Their finite-contour ideas are compatible with the independently detailed
rectangle argument here. The evidence file records exactly what was
inspected; it does not claim an exhaustive audit of the whole repository.

There is a related proof-theoretic pitfall in the literature. Avigad
explains that an appeal to conservation of $I\Sigma_1$ over PRA does not
by itself establish the usual quantifier-alternating statement of the
PNT \cite[footnote 14]{Avigad}. Our arithmetic target has a prescribed
modulus and a quantifier-free matrix, and our argument describes its
finite proof directly. We make no blanket conservation claim for arbitrary
analytic theorems.

\section{Reproduction and verification record}

The archive contains the complete LaTeX sources, the compiled PDF,
the finite rational algorithms, Lean outlines, source provenance, and a
machine-readable verification record. The Python checks verify finite
Abel identities exactly, check the cleared-denominator matrix against
the equivalent rational inequality, and compare a small set of logarithm
and Li enclosures with an independent 100-digit calculation. They also
check representative integer prefactor inequalities. These are useful
implementation checks; none is represented as a computational proof of
the universal PNT or of the zero-free strip.

Rebuild the document with two runs of \texttt{pdflatex} from the package
directory. Run \texttt{python3 code/check\_arithmetic.py} for the finite
checks; this requires \texttt{mpmath} only for the independent numerical
comparisons. The arithmetic specification itself uses Python's standard
library. The final PDF was rendered for layout inspection. Detailed
status and commands are in \texttt{README.txt}.


\begin{thebibliography}{9}
\raggedright
\bibitem{ProveIt}
Vladimir Reshetnikov, \emph{ProveIt}, especially the
\texttt{ConciseIPRA} development,
revision \texttt{b5a3ad037723c44d69a3d5b28a950c554327a989}.
\href{https://github.com/VladimirReshetnikov/ProveIt/tree/b5a3ad037723c44d69a3d5b28a950c554327a989/Logic/PrimitiveRecursiveArithmetic/ConciseIPRA}{Pinned ConciseIPRA source tree}.
Relevant source paths and blob hashes appear in the archive's evidence manifest.

\bibitem{Kontorovich}
Alex Kontorovich and contributors, \emph{Prime Number Theorem And\ldots},
project blueprint, inspected November 15, 2026, especially the second
analytic approach and the $\log^{-9}$ zero-free region.
\url{https://alexkontorovich.github.io/PrimeNumberTheoremAnd/blueprint.pdf}.

\bibitem{Tao}
Terence Tao, \emph{254A, Notes 2: Complex-analytic multiplicative number theory},
December 9, 2014.
\href{https://terrytao.wordpress.com/2014/12/09/254a-notes-2-complex-analytic-multiplicative-number-theory/}{Lecture notes on the author's website}.

\bibitem{Trudgian}
Tim Trudgian, \emph{Updating the error term in the prime number theorem},
arXiv:1401.2689, version 2, 2014.
\url{https://arxiv.org/abs/1401.2689}.
Used to distinguish sharper classical numerical estimates from the
finite constructive estimates proved here.

\bibitem{Avigad}
Jeremy Avigad, \emph{Number theory and elementary arithmetic},
Philosophia Mathematica (3) 11 (2003), 257--284, especially footnote 14.
\url{https://www.andrew.cmu.edu/user/avigad/Papers/elementary.pdf}.

\bibitem{Question}
\emph{Prime number theorem: Prime-counting function in terms of the
logarithmic integral}, Wikipedia, the section supplied in the question,
inspected November 15, 2026.
\url{https://en.wikipedia.org/wiki/Prime_number_theorem}.
This identifies the requested comparison family; the proofs above do not
use it as a source of unproved analytic inputs.
\end{thebibliography}
\end{document}
